\Needspace{0.65\textheight}
\section{The test programs}
\label{app:cases}

All objects are requested with \code{st\_alloc(24)} unless stated otherwise. In the Arena and the memsys5 layer a 24-byte object occupies a 32-byte slot, and in the Free list a 32-byte class slot. ``Live neighbour'' means an object allocated next to the first one and not freed. Each directory \code{corpus/cases/caseNN-*/} contains the program, its \code{case.json} with per-tool expectations and reasons, and a README with a layout diagram.

\begin{center}\small
\begin{tabularx}{\textwidth}{rX>{\raggedright\arraybackslash}p{5.2cm}}
\toprule
\# & trigger & negative control \\
\midrule
0 & two objects, then write \code{a[24]}, one byte past the requested size & write \code{a[23]} \\
1 & read \code{a[24]} & read \code{a[23]} \\
2 & write one byte at \code{a[28]}, four bytes beyond the requested boundary, inside alignment slack & write \code{a[23]} \\
3 & allocate \code{a}, \code{b}, \code{c}; write \code{a[64]}, which is \code{c[0]} in the Arena & write \code{a[23]} \\
4 & read \code{a[pool + 8192]}, past the end of the pool wherever \code{a} landed & read \code{a[23]} \\
5 & \code{a} allocated after another object, then write \code{a[-1]} & write \code{a[0]} \\
6 & read \code{a[-1]} & read \code{a[0]} \\
7 & write the 8 bytes at \code{a - 8}, where a header would be & write \code{a[0..7]} \\
8 & free \code{a}, then write \code{a[0]} (no reuse in between) & write before the free \\
9 & free \code{a}, then read \code{a[0]} & read before the free \\
10 & free \code{a}, \code{b = st\_alloc(24)} receives \code{a}'s storage, then write through \code{a} & write through \code{b} \\
11 & \code{b = st\_realloc(a, 4096)} moves the object, then read through \code{a} & read through \code{b} \\
12 & destroy the whole pool (\code{st\_alloc\_fini}), then read \code{a[0]} & read before the destruction \\
13 & free \code{a} twice & free \code{a} and a second object once each \\
14 & free \code{a + 8} & free \code{a} \\
15 & free \code{a}, \code{b} receives its storage, then free \code{a} again, releasing \code{b}'s storage & free \code{b} \\
\bottomrule
\end{tabularx}
\end{center}

Cases 10 and 15 require immediate reuse, case 11 a moving \code{realloc} and case 12 a pool. The programs print \code{ST NOTE reuse=1}, \code{moved=1} and \code{pool=1} when the property holds, and the harness records the row as unsupported when it does not.

\section{Reproducing the results}
\label{app:repro}

From a clean clone on Ubuntu 26.04 x86-64, the server runner builds the four ASan/Memcheck matrices and the allocator microbenchmark: {\small
\begin{verbatim}
bash tools/setup-server.sh
P=results/$(date +%F)-lab
bash tools/run-allocator-experiments.sh "$P"
python3 harness/render-matrix.py $P-* --check
python3 harness/audit-missed.py $P-* --list
python3 harness/run-real.py --tools asan valgrind --out $P-real
\end{verbatim}}
HWASan uses \code{-{}-tools hwasan} in the AArch64 VM described in \code{docs/hwasan.md}. SQLite and Lua have separate benchmark scripts under \code{benchmarks/}. Archived input bundles, tool environments and source revisions are indexed in \code{results/README.md}; the benchmark summary CSVs identify each performance bundle. \code{make -C report tables} regenerates report tables, and \code{python report/scripts/check-repro.py} verifies regeneration of tables and figures from archived results in an isolated copy. This artifact check does not rerun the experimental programs or replace a full clean-clone execution.
